% GENERATED FILE. Edit canonical lessons, then run npm run generate:books.
% canonical-source-hash: fnv1a64:f248830877430544
% canonical-lessons: LA-C81-frumentum, LA-C81-farina, LA-C81-pirum, LA-C81-ficus, LA-C81-nux

\chapter{Grain, Flour, Pear, Fig, Nut}
\label{ch:la-grain-flour-pear-fig-nut}
\hlchaptermodality{81}

\begin{chapteropening}
\textbf{By the end of this chapter:} \emph{I can name grain, flour, a pear, a fig and a nut.}

Complete the last lesson of chapter 81: i can name grain, flour, a pear, a fig and a nut.
\end{chapteropening}

\section[frūmentum, frūmentī]{frūmentum, frūmentī — grain, corn}

\label{lesson:LA-C81-frumentum}



\begin{quote}
Before the new one: say the Latin for vegetables, then the Latin for an olive.
\end{quote}

\subsection*{You'll want to know: frūmentum, frūmentī}
\textbf{frūmentum, frūmentī} — \textquotedblleft{}grain, corn\textquotedblright{}.

\begin{grammarlens}[title={What you've built}]
One more word for this chapter.
\end{grammarlens}

\subsection*{Your turn}
\begin{itemize}
\raggedright
  \item \emph{Say it:} \emph{frūmentum}
  \item \emph{Say it:} \emph{frūmentum}, once more
  \item \emph{Say it:} say \emph{frūmentum}
  \item \emph{Recall:} say the Latin for vegetables, then the Latin for an olive, then say \emph{frūmentum} again
\end{itemize}

\begin{tcolorbox}[breakable,colback=teal!4,colframe=teal!35!black,title={Before you move on}]
What does frūmentum, frūmentī mean? (\textquotedblleft{}Grain, corn\textquotedblright{}.) Say it once more.
\end{tcolorbox}
\section[farīna, farīnae]{farīna, farīnae — flour}

\label{lesson:LA-C81-farina}



\begin{quote}
Before the new one: say the Latin for an olive, then the Latin for grain.
\end{quote}

\subsection*{You'll want to know: farīna, farīnae}
\textbf{farīna, farīnae} — \textquotedblleft{}flour\textquotedblright{}. English \emph{farinaceous}.

\begin{grammarlens}[title={What you've built}]
One more word for this chapter.
\end{grammarlens}

\subsection*{Your turn}
\begin{itemize}
\raggedright
  \item \emph{Say it:} \emph{farīna}
  \item \emph{Say it:} \emph{farīna}, once more
  \item \emph{Say it:} say \emph{farīna}
  \item \emph{Recall:} say the Latin for an olive, then the Latin for grain, then say \emph{farīna} again
\end{itemize}

\begin{tcolorbox}[breakable,colback=teal!4,colframe=teal!35!black,title={Before you move on}]
What does farīna, farīnae mean? (\textquotedblleft{}Flour\textquotedblright{}.) Say it once more.
\end{tcolorbox}
\section[pirum, pirī]{pirum, pirī — a pear}

\label{lesson:LA-C81-pirum}



\begin{quote}
Before the new one: say the Latin for grain, then the Latin for flour.
\end{quote}

\subsection*{You'll want to know: pirum, pirī}
\textbf{pirum, pirī} — \textquotedblleft{}a pear\textquotedblright{}.

\begin{grammarlens}[title={What you've built}]
One more word for this chapter.
\end{grammarlens}

\subsection*{Your turn}
\begin{itemize}
\raggedright
  \item \emph{Say it:} \emph{pirum}
  \item \emph{Say it:} \emph{pirum}, once more
  \item \emph{Say it:} say \emph{pirum}
  \item \emph{Recall:} say the Latin for grain, then the Latin for flour, then say \emph{pirum} again
\end{itemize}

\begin{tcolorbox}[breakable,colback=teal!4,colframe=teal!35!black,title={Before you move on}]
What does pirum, pirī mean? (\textquotedblleft{}A pear\textquotedblright{}.) Say it once more.
\end{tcolorbox}
\section[fīcus, fīcī]{fīcus, fīcī — a fig, a fig tree}

\label{lesson:LA-C81-ficus}



\begin{quote}
Before the new one: say the Latin for flour, then the Latin for a pear.
\end{quote}

\subsection*{You'll want to know: fīcus, fīcī}
\textbf{fīcus, fīcī} — \textquotedblleft{}a fig, a fig tree\textquotedblright{}.

\begin{grammarlens}[title={What you've built}]
One more word for this chapter.
\end{grammarlens}

\subsection*{Your turn}
\begin{itemize}
\raggedright
  \item \emph{Say it:} \emph{fīcus}
  \item \emph{Say it:} \emph{fīcus}, once more
  \item \emph{Say it:} say \emph{fīcus}
  \item \emph{Recall:} say the Latin for flour, then the Latin for a pear, then say \emph{fīcus} again
\end{itemize}

\begin{tcolorbox}[breakable,colback=teal!4,colframe=teal!35!black,title={Before you move on}]
What does fīcus, fīcī mean? (\textquotedblleft{}A fig, a fig tree\textquotedblright{}.) Say it once more.
\end{tcolorbox}
\section[nux, nucis]{nux, nucis — a nut}

\label{lesson:LA-C81-nux}



\begin{quote}
Before the new one: say the Latin for a pear, then the Latin for a fig.
\end{quote}

\subsection*{You'll want to know: nux, nucis}
\textbf{nux, nucis} — \textquotedblleft{}a nut\textquotedblright{}.

\begin{grammarlens}[title={What you've built}]
One more word for this chapter.
\end{grammarlens}

\subsection*{Your turn}
\begin{itemize}
\raggedright
  \item \emph{Say it:} \emph{nux}
  \item \emph{Say it:} \emph{nux}, once more
  \item \emph{Say it:} say \emph{nux}
  \item \emph{Recall:} say the Latin for a pear, then the Latin for a fig, then say \emph{nux} again
\end{itemize}

\begin{tcolorbox}[breakable,colback=teal!4,colframe=teal!35!black,title={Before you move on}]
What does nux, nucis mean? (\textquotedblleft{}A nut\textquotedblright{}.) Say it once more.
\end{tcolorbox}
